\chapter{Segurança cibernética em automação industrial}
\label{cap:seguranca}

\begin{objetivos}
Ao final deste capítulo você deve ser capaz de:
\begin{itemize}
  \item distinguir segurança funcional (\emph{safety}) de segurança cibernética (\emph{security});
  \item aplicar os conceitos de zona, conduto e nível de segurança da série IEC 62443;
  \item reconhecer a superfície de ataque de um controlador industrial;
  \item escolher contramedidas proporcionais ao risco de cada zona;
  \item descrever um plano de resposta a incidente sem improviso.
\end{itemize}
\end{objetivos}

\section{Safety e security: duas palavras, dois problemas}

O português usa ``segurança'' para as duas coisas, e essa coincidência esconde uma diferença que
decide projetos.

\begin{table}[H]
  \centering
  \caption{Segurança funcional e segurança cibernética.}
  \label{tab:safety-security}
  \begin{tabular}{p{2.6cm} p{5.6cm} p{6.1cm}}
    \toprule
    Aspecto & Segurança funcional (\emph{safety}) & Segurança cibernética (\emph{security}) \\
    \midrule
    Protege contra
      & Falha do sistema que cause dano a pessoas, ao ambiente ou ao patrimônio
      & Ação humana intencional que explore uma vulnerabilidade \\
    \addlinespace
    Pergunta central
      & O sistema falha de forma segura?
      & Alguém de fora consegue fazer o sistema fazer o que não deveria? \\
    \addlinespace
    Base normativa usual
      & Séries IEC 61508 (genérica) e IEC 61511 (indústria de processo), com níveis de integridade de
        segurança (SIL)
      & Série IEC 62443, para sistemas de automação e controle; ISO/IEC 27001 para o sistema de gestão
        corporativo \\
    \addlinespace
    Métrica típica
      & Probabilidade de falha perigosa; nível SIL
      & Nível de segurança (SL) por zona, e requisitos atendidos \\
    \addlinespace
    Documento típico
      & Análise de risco e matriz de causa e efeito dos intertravamentos
      & Análise de risco por zona, arquitetura de rede e plano de resposta \\
    \bottomrule
  \end{tabular}
\end{table}

As duas se cruzam em um ponto perigoso, e a Tabela~\ref{tab:safety-security} resume a distinção
antes de detalhá-la: \textbf{um incidente cibernético pode anular a proteção funcional}. Se um invasor altera a lógica que desliga a máquina, o sistema perde a característica de
segurança que havia sido projetada e validada. É por isso que proteção de máquina não se resolve só
com projeto de \emph{safety}: ela depende também de quem consegue escrever no controlador.

\index{segurança funcional}\index{segurança cibernética}

\section{A base normativa: a série IEC 62443}

A série \textbf{IEC 62443} trata da segurança cibernética de sistemas de automação e controle
industriais. Ela nasceu do comitê ISA99, em 2002, e passou a ser desenvolvida em conjunto com o grupo
de trabalho IEC TC65 WG10; hoje é publicada pela ISA como ANSI/ISA 62443 nos Estados Unidos e pela
IEC como IEC 62443 internacionalmente, com conteúdo técnico equivalente quando as duas organizações
aceitam o texto. A IEC a reconhece como norma horizontal: normas setoriais de segurança em
tecnologia de operação devem usá-la como base.

Os conceitos que interessam ao profissional de automação:

\begin{itemize}
  \item \textbf{papéis}: o proprietário do ativo, o fornecedor do produto e o prestador de serviço
        têm responsabilidades diferentes. ``Segurança cibernética é responsabilidade do fabricante''
        é uma frase falsa em qualquer um dos três casos;
  \item \textbf{zonas e condutos}: a instalação é dividida em zonas --- agrupamentos de ativos com
        requisitos de segurança comuns --- e os condutos são os caminhos de comunicação autorizados
        entre zonas. É o conceito que substitui a ideia de ``rede plana'', em que tudo fala com tudo;
  \item \textbf{níveis de segurança (SL 0 a 4)}: indicam a resistência exigida contra classes
        distintas de atacante. Na definição da série, o nível 0 dispensa proteção; o nível 1 protege
        contra uso indevido não intencional; o nível 2, contra uso indevido intencional por meios
        simples e poucos recursos; o nível 3, contra meios sofisticados e recursos moderados; e o
        nível 4, contra meios sofisticados com recursos extensos. A avaliação é feita requisito por
        requisito, e não como ``nota'' do produto;
  \item \textbf{níveis de maturidade (1 a 4)}: avaliam o processo de quem \emph{desenvolve} o produto
        ou o serviço, do trabalho informal até a melhoria contínua medida.
\end{itemize}

\begin{figure}[H]
  \centering
  \includegraphics[width=0.98\linewidth]{22-zonas-condutos.pdf}
  \caption{Zonas com requisitos distintos e os condutos que autorizam a comunicação entre elas.}
  \label{fig:zonas-condutos}
\end{figure}

A série é revisada continuamente e publicada em partes. Para citar uma parte é preciso citar a edição
e o ano, e conferir a versão vigente na fonte --- por exemplo, entre as partes publicadas estão
IEC 62443-3-3 (requisitos de segurança de sistema, 2013), IEC 62443-4-1 (ciclo de vida de
desenvolvimento seguro, 2018), IEC 62443-4-2 (requisitos técnicos de componentes, 2019),
IEC 62443-3-2 (avaliação de risco para projeto de sistema, 2020) e IEC 62443-2-1 (programa de
segurança do sistema de automação, 2024). Como a série evolui, o ano citado aqui corresponde à
publicação consultada, não a uma promessa de permanência.

Para o lado corporativo e de gestão, dois documentos complementam o quadro:

\begin{itemize}
  \item \textbf{ISO/IEC 27001:2022} (terceira edição, com emenda de 2024): requisitos para um sistema
        de gestão de segurança da informação. Cobre a empresa inteira, inclusive a parte
        administrativa; não substitui a IEC 62443 no que é específico de controle industrial;
  \item \textbf{NIST Cybersecurity Framework, versão 2.0}: referência de programa, organizada em
        funções de governança, identificação, proteção, detecção, resposta e recuperação. Serve como
        roteiro de maturidade, não como especificação técnica.
\end{itemize}

\index{IEC 62443}\index{zonas e condutos}\index{ISO/IEC 27001}

\section{A superfície de ataque do controlador}

Um controlador industrial não é um computador comum, mas compartilha com ele uma característica:
possui vários pontos de entrada. A Tabela~\ref{tab:superficie-ataque} lista os que aparecem com mais
frequência e a contramedida correspondente.

\begin{table}[H]
  \centering
  \caption{Pontos de entrada de um sistema de automação e contramedidas usuais.}
  \label{tab:superficie-ataque}
  \begin{tabular}{p{3.5cm} p{5.3cm} p{5.5cm}}
    \toprule
    Ponto de entrada & Como o problema aparece & Contramedida \\
    \midrule
    Ambiente de engenharia
      & Estação de programação com acesso à rede de controle e à internet
      & Estação dedicada, sem navegação, em zona própria; acesso apenas pelo conduto autorizado \\
    \addlinespace
    Serviços de rede expostos
      & Servidor web do controlador, FTP, telnet ou diagnóstico abertos
      & Desabilitar o que não é usado; inventariar portas e serviços por equipamento \\
    \addlinespace
    Credenciais padrão
      & Senha de fábrica mantida em painel, IHM ou gateway
      & Trocar antes do comissionamento; conta por pessoa, não compartilhada \\
    \addlinespace
    Mídia removível
      & Pen drive para levar o programa até a máquina
      & Política de mídia controlada e dedicada; verificação antes do uso \\
    \addlinespace
    \emph{Firmware} e atualização
      & Versão antiga, sem correção, ou atualização sem controle
      & Inventário de versões e janela de atualização planejada \\
    \addlinespace
    Rede plana
      & Rede de controle no mesmo domínio da rede administrativa
      & Segmentação em zonas, com conduto intermediado por \emph{firewall} industrial \\
    \addlinespace
    Acesso remoto
      & Acesso direto para suporte do fornecedor, sem registro
      & Acesso por VPN com autenticação, sessão registrada e autorizada por chamado \\
    \bottomrule
  \end{tabular}
\end{table}

\section{Casos públicos que vale conhecer}

Estudar incidentes reais é mais eficiente que decorar taxonomia. Três exemplos são públicos,
documentados e suficientes para fixar o problema --- e nenhum deles precisa ser detalhado em termos de
exploração:

\begin{itemize}
  \item \textbf{Stuxnet (2010)}. Programa malicioso que atingiu sistemas de controle industrial e
        demonstrou publicamente que uma instalação física pode ser alvo de manipulação de lógica. É o
        caso que tirou a segurança cibernética industrial da pauta teórica;
  \item \textbf{vulnerabilidade de credencial em controladores Logix (2021, CVE-2021-22681)}. O
        ambiente de engenharia e as versões de \emph{firmware} correspondentes usavam uma chave para
        verificar se quem se comunica com o controlador é o software de programação legítimo. Um
        atacante não autenticado poderia burlar essa verificação e autenticar-se no controlador. A
        severidade é crítica (base 9,8 na escala CVSS 3.1 do NVD) e o caso consta do catálogo de
        vulnerabilidades exploradas conhecidas da agência norte-americana CISA;
  \item \textbf{injeção de código pelo ambiente de engenharia (2022, CVE-2022-1159)}. Um atacante que
        obtivesse acesso de administrador na estação de trabalho com o software de programação poderia
        inserir código no controlador de forma indetectável ao usuário. Severidade alta (base 7,2 no
        NVD e 7,7 na avaliação do ICS-CERT). O caso mostra que o ponto fraco não é o controlador, é a
        \textbf{estação de engenharia}.
\end{itemize}

\begin{atencao}
O padrão dos três casos é o mesmo: o alvo preferido não é a lógica de controle, é o \emph{caminho}
até ela --- a estação de engenharia, a credencial, o serviço aberto. Investir em proteção do
controlador e deixar a estação de programação navegando na internet no mesmo segmento de rede é
proteger a porta e deixar a janela aberta.
\end{atencao}

\section{Contramedidas, do projeto à operação}

A ordem abaixo segue o princípio de que arquitetura vale mais que produto: nenhum equipamento resolve
o que a segmentação resolve.

\begin{enumerate}
  \item \textbf{segmentar em zonas e definir os condutos}, como na Figura~\ref{fig:zonas-condutos}. É a
        decisão de maior efeito e a mais barata quando tomada no projeto;
  \item \textbf{intermediar todo tráfego entre zonas} por \emph{firewall} industrial, com regras
        explícitas por endereço, porta e direção. Sem regra escrita, não há conduto;
  \item \textbf{controlar o acesso remoto}: VPN com autenticação individual, sessão registrada e
        autorização por chamado. Acesso de fornecedor é acesso privilegiado e deve ser tratado como
        tal;
  \item \textbf{eliminar credenciais padrão} e individualizar contas. Nenhum sistema de automação
        deve ter usuário ``admin'' compartilhado por quatro pessoas;
  \item \textbf{desabilitar serviços não utilizados}: o servidor web do controlador que ninguém usa é
        superfície de ataque sem contrapartida;
  \item \textbf{controlar mídia removível}, com mídia da empresa, verificada e destinada a esse uso;
  \item \textbf{fazer \emph{hardening} de IHM e dos computadores industriais}: sistema atualizado,
        serviços mínimos, sem navegação e sem \emph{e-mail};
  \item \textbf{registrar e fazer backup}: registro de eventos do controlador, da rede e dos acessos;
        backup do projeto, guardado fora do equipamento e testado. Backu\emph{p} que nunca foi
        restaurado não é backup, é expectativa;
  \item \textbf{gerenciar atualização com janela planejada}. Em automação, a correção de segurança
        compete com a parada de produção: a decisão é de risco, tomada com inventário de versões em
        mãos;
  \item \textbf{treinar quem opera}. A maior parte dos incidentes começa em uma ação humana ---
        conexão indevida, mídia desconhecida, senha anotada.
\end{enumerate}

\section{Resposta a incidente}

Ter plano de resposta não impede o incidente; impede o improviso durante ele. O plano precisa caber
em uma página e responder a cinco perguntas:

\begin{enumerate}
  \item \textbf{quem decide} desligar ou isolar um trecho da planta, e quem é o substituto;
  \item \textbf{qual é o modo seguro} para a planta operar sem o sistema comprometido --- operação
        manual, com supervisão local, ou parada programada;
  \item \textbf{de onde se restaura}: qual é a última versão confiável do programa e da configuração,
        e onde ela está guardada;
  \item \textbf{como se registra a evidência}: anotar horário, telas, registro de eventos e o que foi
        alterado, antes de mexer --- inclusive para permitir análise posterior;
  \item \textbf{como se volta à produção}: retorno gradual, com conferência de lógica, intertravamento
        e teste dos casos críticos, como no Capítulo~\ref{cap:simulacao}.
\end{enumerate}

Duas armadilhas comuns nessa hora: restaurar um backup que já contém o código alterado pelo atacante
--- daí a importância de guardar versões fora da rede --- e reiniciar o equipamento ``para ver se
resolve'', o que apaga evidência sem resolver a causa.

\section{Segurança é processo, não projeto}

Um sistema de automação é seguro por decisões mantidas ao longo do tempo: inventário atualizado,
zonas respeitadas, acessos revisados, versões acompanhadas e equipe treinada. A parte técnica é a
menor parte do trabalho; a maior é disciplina.

No contexto do curso técnico, isso se traduz em três hábitos que valem desde a bancada:

\begin{itemize}
  \item \textbf{proteger desde o primeiro teste}: trocar senha padrão, não deixar o laboratório em
        rede aberta e não usar pen drive de origem desconhecida;
  \item \textbf{documentar o que foi conectado}: qualquer equipamento novo em rede, com endereço,
        porta e finalidade (Capítulo~\ref{cap:documentacao});
  \item \textbf{tratar o acesso como privilégio}: quem programa altera o comportamento físico de uma
        máquina. Esse poder não deve estar disponível em qualquer computador.
\end{itemize}

\begin{pratica}
Desenhe as zonas e os condutos de uma planta que você conheça, no formato da
Figura~\ref{fig:zonas-condutos}, e depois:
(a) atribua o nível de segurança que você julgar adequado a cada zona, justificando;
(b) liste os serviços que estão abertos hoje e diga quais poderiam ser desabilitados;
(c) escreva as três primeiras regras de \emph{firewall} entre as zonas;
(d) redija o plano de resposta a incidente em uma página, respondendo às cinco perguntas da seção
anterior.
\end{pratica}

\begin{exercicios}
\begin{enumerate}[label=\alph*)]
  \item Explique a diferença entre \emph{safety} e \emph{security} usando um exemplo de cada.
  \item Por que um incidente cibernético pode anular a proteção funcional de uma máquina?
  \item Descreva o conceito de zona e o de conduto, e diga por que ele substitui a ideia de ``rede
        plana''.
  \item Explique os níveis de segurança 1 a 4 da série IEC 62443, com suas próprias palavras.
  \item Qual é a diferença entre nível de segurança e nível de maturidade? A quem cada um se aplica?
  \item Cite três papéis definidos pela série IEC 62443 e diga uma responsabilidade de cada um.
  \item Um gerente afirma: ``segurança cibernética da planta é responsabilidade do fabricante do
        controlador''. Responda com base nos papéis definidos pela norma.
  \item Escolha quatro linhas da Tabela~\ref{tab:superficie-ataque} e descreva a contramedida
        correspondente em uma frase cada.
  \item O que os casos CVE-2021-22681 e CVE-2022-1159 têm em comum quanto ao ponto de entrada?
  \item Por que a estação de engenharia costuma ser o alvo preferido? Cite duas medidas práticas de
        proteção.
  \item Explique por que ``backup que nunca foi restaurado não é backup''.
  \item Escreva as cinco perguntas que o plano de resposta a incidente precisa responder.
  \item Por que restaurar um backup antigo pode não resolver um incidente, e como evitar esse
        problema?
  \item Escolha uma medida da seção de contramedidas que dependa de disciplina, e não de produto, e
        explique como você garantiria que ela seja mantida na rotina da planta.
\end{enumerate}
\end{exercicios}
